% Methods M07 to M12.

\subsection{Phase-fixed weighted GS (M07)}\label{sec:m07}

Phase-fixed weighted GS is the practical method for large arrays. The solution replicated is that of Kim, Keesling, Omran, Levine, Bernien, Greiner, Lukin and Englund~\cite{kim2019}: after an initial number $N$ of weighted-GS iterations the focal-plane phase is frozen, so that later iterations adjust amplitudes only and the weights stay effective. The same procedure is used in the experiments with a thousand tweezers of Ebadi et al.~\cite{ebadi2021} and, through the references of Manetsch et al.~\cite{manetsch}, at the scale of thousands of sites. The source reports a modulation efficiency of 91.2 percent and a non-uniformity below 0.5 percent on $50\times30$ foci, and a camera correction from 22 percent to 1.59 percent in three steps against 6.47 percent for conventional weighted GS.

\textbf{Derivation.} At iteration $i$ let $B_i(\mathbf u_m)=|V_m|/a_m$ be the focal amplitude and $T_m=a_m$ the target. With $\langle B_i\rangle_N$ the mean over the traps, the weights and the SLM field follow
\begin{equation}
g_i(\mathbf u_m)=\frac{\langle B_i\rangle_N}{B_i(\mathbf u_m)}\,g_{i-1}(\mathbf u_m),\quad g_0=1,\qquad
\mathcal F\big[A_{i+1}e^{i\phi_{i+1}}\big]=g_i\,T\,e^{i\psi_i},
\label{eq:m07}
\end{equation}
where the focal phase follows the field, $\psi_i=\arg V_m$, for the first $N$ iterations and is then frozen at its value $\psi_N$,
\begin{equation}
\psi_i=\psi_N\quad(i\ge N).
\label{eq:m07-freeze}
\end{equation}
In the engine the first iteration projects the field of the start hologram with $g=1$ and the first weight update uses the field of the second iteration, so that Eq.~\eqref{eq:m07} holds from $i\ge2$ with $g_1=g_0$. In the recursion of Kim et al.~\cite{kim2019} the weight $g_i$, built from the focal amplitude $B_i$ of the $i$-th iteration, produces the hologram of the $(i+1)$-th, with $g_0=1$ (preprint arXiv:1903.09286, page 2). The text read does not say whether $g_0=1$ is applied to the field of the start hologram, as the engine does, or whether the update begins at the first measured field, which would be one iteration earlier than the engine; the preprint text available to this report does not settle the indexing of the first iteration. The engine follows the phase-fixed weighted GS of the open library slmsuite (no update at its iteration 0); the convention is DinHot's choice and not a verified statement of the source. The weights $g$ keep updating after the freeze and carry the memory of all earlier corrections. The reason given by the source is that replacing the SLM-plane amplitude by the laser amplitude introduces a change $\delta\psi_i$ of the focal phase, which makes $g_i$ less effective; fixing $\psi$ removes the change. The sources do not define $B_i$ (it is taken as the field amplitude, as in weighted GS) nor say whether the $N$-th iteration already uses the frozen phase; here it computes the phase, and a one-iteration shift turns the parity test of the port red, which is how the choice is pinned. The source states the weight update and the freezing of the focal phase in its section on the calculation of the hologram and in its Fig.~1 (Kim et al.~\cite{kim2019}); both are displayed without an equation number in the preprint, so they are cited here by section and figure.

\textbf{Pseudocode.} Algorithm~\ref{alg:m07} lists the procedure as implemented in the Fourier form.

\begin{algorithm}[t]
\caption{Phase-fixed weighted GS in the Fourier form. The weights $g$ are updated at every iteration but the first, which projects the field of the start hologram with $g=1$; the focal phase $\psi$ is recomputed until iteration $N$ and then kept.}
\label{alg:m07}
\begin{algorithmic}[1]
\Require traps $\{\mathbf u_m\}$ on FFT bins, target amplitudes $a_m$, freeze index $N$, iteration budget $K$, start $\phi^0$
\State $g_m\gets 1$ for all $m$; $\phi\gets\phi^0$ (random superposition when no start is given)
\For{$k=1,\dots,K$}
  \State $V\gets\mathcal F[A e^{i\phi}]$ sampled at the trap bins; $b_m\gets|V_m|/a_m$
  \If{$k>1$} \State $g_m\gets g_m\,\langle b\rangle/b_m$ \EndIf
  \If{$k<N$} \State $\psi_m\gets\arg V_m$ \ElsIf{$k=N$} \State $\psi_m\gets\arg V_m$; store $\psi$ \EndIf
  \State $c_m\gets g_m\,a_m\,e^{i\psi_m}$ (with the stored $\psi$ for $k>N$)
  \State $\phi\gets\arg\mathcal F^{-1}\big[\sum_m c_m\delta(\mathbf u-\mathbf u_m)\big]$
\EndFor
\State \Return $\phi$
\end{algorithmic}
\end{algorithm}

\textbf{Parameters and how they were chosen.} The freeze index has the default \defv{M07}{nfix}, the value of Kim et al., pinned by a test; the acceptance criterion is run with $N=60$, which was measured on the $50\times30$ task of DinHot's grid, and the default is not grid dependent (whether it should become grid dependent is not settled). The measured finding is that the literal setting does not transfer to a smaller SLM: with the source value the spread stalls on the crowded array, and larger values of $N$ remove the stall (Section~\ref{sec:results}). A replication of a published setting that does not transfer is itself a result and is kept. Since the corrections of the update rule, the method matches the weighted GS with the phase freeze of the open library slmsuite to a small fraction of a percent in the final spread from the same start; this holds with the convention above, in which the first iteration is plain GS; the engine before that convention reweighted the field of the start hologram.

\textbf{Complexity.} $O(P\log P)$ per iteration in the Fourier form with $P=n_{\mathrm{fft}}^2$ and $O(N_{\mathrm{pix}}M)$ in the point form; the method is two-dimensional only because the freeze is defined for a plane. The WebGPU backend covers this method with the traps on bins and the ideal device (Section~\ref{sec:live}).

\textbf{Acceptance criterion and gate.} The ladder row requires fewer iterations than weighted GS to reach a relative standard deviation of 0.5 percent on a $50\times30$ array with the phase fixed after the measured index (Table~\ref{tab:ladder}). The gate runs two seeds of 150 iterations in the Fourier form and asserts that phase-fixed weighted GS reaches the target and that weighted GS either never does or needs more iterations, with a final efficiency above 0.6; a second test pins the stall with the source value. The measurement of this report is: \acc{M07}. A sweep of the freeze index is also run by the script of this report (the key \texttt{n\_fix\_sweep\_min\_rel\_std\_150} of the acceptance file; Table~\ref{tab:acceptance}).

\textbf{Limits.} The index depends on the grid; on a device with crosstalk and four phase levels the method keeps the spread of the ideal-model optimum, which is the gap that the first proposed method removes (Section~\ref{sec:b1}). The module is \texttt{m07\_phase\_fixed\_wgs}.

\subsection{Mixed-region amplitude freedom (M08)}\label{sec:m08}

Mixed-region amplitude freedom (MRAF) is the method for smooth potentials, a disc, a ring or a bar, in which GS leaves speckle. The solution replicated is that of Pasienski and DeMarco~\cite{pasienski}: the amplitude is left free outside the signal region, so that the power the shape does not need drains into a noise region and the shape comes out accurate. The same rule, with a mixing parameter $p$ in place of $m$, is Eq.~(3) of Holland et al.~\cite{holland}, who use $p=0.7$ for a lattice of traps. With signal region SR, noise region NR and mixing parameter $m$ the focal constraint is
\begin{equation}
G^{(n)}=\Big\{m\,T\big|_{\mathrm{SR}}+(1-m)\,\big|E^{(n)}_{\mathrm{out}}\big|\,\Big|_{\mathrm{NR}}\Big\}\,e^{\,i\arg E^{(n)}_{\mathrm{out}}},\qquad
\phi^{(n+1)}=\arg\mathcal F^{-1}G^{(n)},
\label{eq:m08}
\end{equation}
with the SLM amplitude reset to the beam. In the signal region the amplitude is forced to $mT$, and in the noise region it is kept and scaled by $1-m$; $m=1$ is plain GS with a zero noise region, which a test pins. The target $T$ is rescaled to the input power. The error of the history is the root-mean-square fractional error of the source,
\begin{equation}
\eta=\sqrt{\Big\langle\Big(\frac{s\,I-T^2}{T^2}\Big)^2\Big\rangle},\qquad s=\frac{\sum T^2}{\sum I},
\label{eq:m08-err}
\end{equation}
over the measure region. Three design decisions are DinHot's and not the source's. First, the measure region is the shape eroded by one resolution cell, because a pixel closer to the edge than the cell cannot reach its target for any hologram; without it the error floor is dominated by the edge, and a shape thinner than the erosion is refused. Second, the default start is a Fresnel-lens phase of \defv{M08}{startwaves} waves at the aperture edge, because a random start leaves speckle that the iteration removes slowly. Third, the offset variants of the method are not implemented, because their primary papers were not opened.

\textbf{Pseudocode.} Algorithm~\ref{alg:m08} lists the iteration; with $m=1$ it is plain GS with a zero noise region.

\begin{algorithm}[t]
\caption{Mixed-region amplitude freedom (the field loop of the engine). The target is rescaled to the input power, the noise region keeps the amplitude it reaches, scaled by $1-m$.}
\label{alg:m08}
\begin{algorithmic}[1]
\Require target amplitude $T$ and signal mask SR on the focal grid, mixing $m\in(0,1]$, start waves $W$, iteration budget $K$, start $\phi^0$
\State $T\gets T\sqrt{N_{\mathrm{pix}}/\sum_{\mathrm{SR}}T^2}$ on SR, $0$ elsewhere; measure region $\gets$ SR eroded by one resolution cell
\State $\phi\gets\phi^0$, or the Fresnel-lens phase of $W$ waves at the aperture edge when no start is given; $E\gets\mathcal F[Ae^{i\phi}]$
\For{$k=1,\dots,K$}
  \State $G\gets m\,T$ on SR, $(1-m)\,|E|$ elsewhere; $G\gets G\,e^{i\arg E}$
  \State $\phi\gets\arg\mathcal F^{-1}[G]$; $E\gets\mathcal F[Ae^{i\phi}]$ (the SLM amplitude is reset to the beam)
  \State record the efficiency in SR and the error $\eta$ of Eq.~\eqref{eq:m08-err} over the measure region
\EndFor
\State \Return $\phi$
\end{algorithmic}
\end{algorithm}

\textbf{Parameters.} The mixing has the default \defv{M08}{mixing}, which lies in the range the source reports as typical; on DinHot's grid a larger value trades a higher efficiency for a larger error, a trade measured when the default was set. The benchmark tunes the mixing over a grid of four values.

\textbf{Complexity.} Two Fourier transforms per iteration, $O(P\log P)$; the source reports convergence in under a hundred iterations.

\textbf{Acceptance criterion and gate.} On a disc, a ring and a bar, the signal-region error is at least five times below GS at the cost of efficiency (Table~\ref{tab:ladder}). The gate runs the three shapes on the live grid with two starts of different focal curvature (the default start is deterministic, so the starts are varied instead of seeds) and asserts a ratio of errors of at least five and a lower efficiency for both. The measurement of this report is: \acc{M08}. A thin line is resolution limited: for a line a few bins wide the ratio falls below five because the error floor is the shape edge and not the method, so the criterion is met for shapes wider than the resolution cell, and the thin-line limit is a recorded result.

\textbf{Why the noise region helps.} In plain GS the whole focal plane outside the traps is constrained to zero amplitude, which forces the field to satisfy a constraint it cannot meet with a phase-only hologram, and the unsatisfiable part reappears as speckle inside the signal region. Freeing the amplitude outside the signal region lets the power that the shape does not need leave through the noise region; the mixing $m$ sets how much of the signal amplitude is imposed against how much of the achieved field is kept in the noise region, so that a smaller $m$ gives the noise region more freedom, a lower efficiency in the signal region and a more accurate shape. The efficiency of the method is therefore the fraction of the power that the iteration leaves in the signal region at its fixed point, which is the quantity that the history reports. A shape thinner than the resolution cell cannot be improved by any choice of $m$, which is the limit that the thin-line test pins.

\textbf{Limits.} The efficiency is low by design. The model has no device: a hologram optimized on the ideal device and scored on the crosstalk device of case C04 shows a signal error that differs by a large factor from the acceptance value, and the single factor responsible is not established, since the case differs from the test in disc radius, margin and device at once; no single-factor comparison was run. The module is \texttt{m08\_mraf} and the field loop of \texttt{\_gs}.

\subsection{Direct search (M09)}\label{sec:m09}

Direct search is the optimization method of holographic optical tweezers: visit every SLM pixel in random order, try every phase level and keep the best against an explicit cost that trades efficiency for uniformity with one weight. It is the reference method for perfect uniformity in Di Leonardo et al.~\cite{dileonardo} and the explicit-cost ancestor of the gradient methods. The cost is that of Polin et al.~\cite{polin}. With $I_m$ the trap intensities as fractions of the input power, $I^D_m=w_m/\operatorname{mean}(w)$ the design intensities and $M$ the trap count,
\begin{equation}
C=-\langle I\rangle+f_u\,\sigma,\qquad \sigma=\sqrt{\frac1M\sum_m\big(I_m-\gamma I^D_m\big)^2},\qquad
\gamma=\frac{\sum_m I_mI^D_m}{\sum_m(I^D_m)^2},
\label{eq:m09}
\end{equation}
where $f_u$ weights uniformity against efficiency (the symbol $f$ is the focal length elsewhere), $\sigma$ is the root-mean-square residual of the design in units of the intensity and not the percentage of Eq.~\eqref{eq:metrics}, and $\gamma$ is the best common scale of the design and not the drag of Eq.~\eqref{eq:bead}. One sweep visits every pixel once in a seeded random order; at each pixel all $P$ trial phases $2\pi k/P$ are evaluated and the best is kept if it lowers the cost, so the cost never increases. The trap fields are updated incrementally: a change $\delta z=e^{i\phi'}-e^{i\phi}$ at pixel $j$ changes each field by
\begin{equation}
\delta V_m=\frac{A_j}{N_{\mathrm{pix}}}\,\delta z\,e^{-i\Delta_j^m},
\label{eq:m09-update}
\end{equation}
an update that costs $O(M)$ and holds for any trap depth and any vortex charge, because it uses the point-trap form. After each sweep the fields are recomputed in full to drop the rounding drift of the increments. The start is the random superposition rounded to the $P$ levels, and an initial phase overrides it. The source reports convergence in about $N$ trials and a worst case of $N\,P$.

\textbf{Pseudocode.} Algorithm~\ref{alg:m09} lists a run.

\begin{algorithm}[t]
\caption{Direct search with the cost of Eq.~\eqref{eq:m09}. One sweep visits every pixel once in a seeded random order; the cost never increases.}
\label{alg:m09}
\begin{algorithmic}[1]
\Require traps with design intensities $I^D_m$, levels $P$, weight $f_u$, number of sweeps $K$, cap on the trial phases
\State $\phi\gets$ the random superposition rounded to the $P$ levels $2\pi p/P$; $V\gets$ trap fields of $\phi$
\For{sweep $=1,\dots,K$}
  \State $C_{\mathrm{cur}}\gets C(|V|^2)$ by Eq.~\eqref{eq:m09}; draw a random order of the $N_{\mathrm{pix}}$ pixels from the seeded generator
  \For{each pixel $j$ in that order}
    \State \textbf{stop} when the cap on the trial phases would be exceeded
    \State for $p=0,\dots,P-1$: $V^{(p)}\gets V+\frac{A_j}{N_{\mathrm{pix}}}\big(e^{2\pi ip/P}-e^{i\phi_j}\big)e^{-i\Delta_j^m}$ and $C^{(p)}\gets C(|V^{(p)}|^2)$
    \State $p^\ast\gets\arg\min_pC^{(p)}$
    \If{$C^{(p^\ast)}<C_{\mathrm{cur}}$} \State $V\gets V^{(p^\ast)}$; $\phi_j\gets2\pi p^\ast/P$; $C_{\mathrm{cur}}\gets C^{(p^\ast)}$ \EndIf
  \EndFor
  \State recompute $V$ in full from $\phi$ (drops the rounding drift of the increments); record the history entry
\EndFor
\State \Return $\phi$
\end{algorithmic}
\end{algorithm}

\textbf{Parameters.} The number of levels has the default \defv{M09}{levels} and the weight $f_u$ the default \defv{M09}{f}, the values of the direct-search row of Di Leonardo et al.; the budget is the number of sweeps taken from the iteration budget, and a cap on the trial phases serves the time-boxed rows of the benchmark. The simulated-annealing variant of the source is not implemented, because the source gives no quantitative result for it.

\textbf{Complexity and lane.} One trial costs $O(M)$ and a sweep evaluates $N_{\mathrm{pix}}P$ trials, so a sweep costs $O(N_{\mathrm{pix}}PM)$; the method is sequential by construction, so it is baked and replayed and not run live. A test confirms that the incremental update agrees with a full evaluation to $10^{-10}$ and that the cost never rises.

\textbf{Acceptance criterion and gate.} The criterion requires a uniformity of at least 0.98 on the $10\times10$ task and a spread below five percent on a quasi-periodic set (Table~\ref{tab:ladder}). The quasi-periodic set is DinHot's own, a product of two Fibonacci chains with a stated gap, with two sites removed to give 119 traps, because the coordinates of the 119 traps of Polin et al. are not published. The gate asserts the uniformity and an efficiency above 0.3 (so that the uniformity is not bought by switching traps off). The measurement of this report is: \acc{M09}. The published reference is a uniformity of 1.00 and an efficiency of 0.68 for direct search on the $10\times10$ task, and a spread above fifty percent for GS and the adaptive form on the quasi-periodic set~\cite{polin,dileonardo}, which the product does not re-run.

\textbf{Why the search can stop.} At every pixel the best of the $P$ trial phases is kept only if it lowers the cost, and the current phase is among the candidates, so each accepted step is a non-increasing step of $C$ and the sequence of sweep costs is monotone and bounded below by the minimum of Eq.~\eqref{eq:m09}. The search therefore converges to a coordinate-wise local minimum, a phase pattern in which no single pixel can be improved by any of the $P$ levels, and not to a global minimum; this is the local minimum that Polin et al. report. The scale $\gamma$ in the definition of $\sigma$ makes the uniformity term insensitive to the overall efficiency, so that the weight $f$ trades the two quantities and does not merely rescale one of them. The result is sensitive to the order of the visits, which the seeded permutation fixes, and to the number of levels, since a lower number of levels leaves each pixel a coarser choice.

\textbf{Limits.} The method is sequential, slow and stuck in local minima; the ghost traps absorb efficiency, so the efficiency in the bake is lower than that of weighted GS. The cost model is the ideal device: scored on a device with crosstalk and four phase levels the uniformity drops, which is the reason for the device-aware methods. The module is \texttt{m09\_direct\_search}.

\subsection{Conjugate-gradient minimization (M10)}\label{sec:m10}

Conjugate-gradient minimization optimizes all pixel phases at once against an intensity cost with a closed-form gradient. The solution replicated is that of Harte, Bruce, Keeling and Cassettari~\cite{harte}, the reference for smooth potentials with suppressed optical vortices and the parent of the device-aware variant of Section~\ref{sec:m13} and of the gradient family. The SLM field is $E_{\mathrm{in}}=A_je^{i\phi_j}$, the focal field $\tilde\psi$ is its zero-padded Fourier transform, and the cost over the signal region (the noise region is free) is
\begin{equation}
C_t=\sum_{nm}\big(T_{nm}-|\tilde\psi_{nm}|^2\big)^t,\qquad
\frac{\partial C_t}{\partial\phi_{rs}}=2t\,\operatorname{Re}\big(i\,\psi_{rs}^*X_{rs}\big),
\label{eq:m10}
\end{equation}
\begin{equation}
X_{rs}=\sum_{nm}(T_{nm}-I_{nm})^{t-1}\,\tilde\psi_{nm}^*\,e^{+2\pi i(rn+sm)/N}.
\label{eq:m10-x}
\end{equation}
In the code the cost is normalized by $\sum(sT)^t$, which has the same minimizers, with the target intensity carrying the requested efficiency $\eta$: for traps $T_m=\eta\,w_m/\sum w$, and for fields $T=\eta\,N_{\mathrm{pix}}a^2/\sum_{\mathrm{SR}}a^2$, so that $\sum T$ is the requested fraction of the total power. The gradient costs two Fourier transforms, one forward and one adjoint, which makes the method cheap per step and lets the TypeScript port run without automatic differentiation. The basic cost $t=2$ leaves optical vortices in the signal region; the powered cost $t=4$ suppresses them, and the sequential $t=4$ then $t=2$ reaches a small error on a ring and power-law target in the source after thousands of iterations.

\textbf{Optimizer.} The Fletcher--Reeves nonlinear conjugate gradient builds directions $d_i=-g_i+\gamma_id_{i-1}$ with $\gamma_i=g_i\cdot g_i/(g_{i-1}\cdot g_{i-1})$ (Harte Eqs.~(3) to (5)), with a strong-Wolfe line search by cubic interpolation. The source does not state the line search; this is DinHot's choice, and a second Wolfe constant below one half keeps the Fletcher--Reeves directions descent directions. If a direction is not a descent direction or the search fails, the iteration restarts from steepest descent and the run stops if that fails too, so that the cost never increases. The start is $n_{\mathrm{pix}}^2$ uniform random phases from the portable generator.

\textbf{Pseudocode.} Algorithm~\ref{alg:m10} lists the optimizer.

\begin{algorithm}[t]
\caption{Fletcher--Reeves conjugate gradient with a strong-Wolfe line search on the powered intensity cost, in stages of equal share of the iterations (the loop shared with the device-aware conjugate gradient, which differs by the forward model).}
\label{alg:m10}
\begin{algorithmic}[1]
\Require target (traps or shaped field), power schedule $(t_1,\dots,t_S)$, requested efficiency $\eta$, iteration budget $K$, start $\phi^0$
\State $\phi\gets\phi^0$ (uniform random phases when no start is given)
\For{each stage $s=1,\dots,S$ with $K_s\approx K/S$ iterations (the last stage takes the remainder)}
  \State build the cost $C_{t_s}$ with target level $\eta$; $f\gets C_{t_s}(\phi)$; $g\gets\nabla C_{t_s}(\phi)$ by Eq.~\eqref{eq:m10} (two Fourier transforms); $d\gets-g$
  \For{$i=1,\dots,K_s$}
    \If{$d\cdot g\ge0$} \State $d\gets-g$ \EndIf
    \State strong-Wolfe search along $d$ ($c_1=10^{-4}$, $c_2=0.1$, cubic interpolation, first trial moving the largest pixel by 0.1 rad, later trials $2(f-f_{\mathrm{old}})/(g\cdot d)$)
    \If{no step is found} \State restart with $d\gets-g$ and search again; \textbf{stop} the run if that fails too \EndIf
    \State $\phi\gets\phi+a\,d$; $g'\gets\nabla C_{t_s}(\phi)$; $\gamma\gets(g'\cdot g')/(g\cdot g)$; $d\gets-g'+\gamma d$; $g\gets g'$
  \EndFor
\EndFor
\State \Return $\phi$ wrapped to $[0,2\pi)$
\end{algorithmic}
\end{algorithm}

\textbf{Parameters.} The power $t$ has the default stages of the field case (the powered cost first, then the basic one, with an equal share of the iterations) and the basic cost for traps; the design-efficiency argument is left unset by default and the code then chooses a design level for traps and a lower one for fields (Table~\ref{tab:defaults}); the Wolfe constants and the cap on line-search evaluations are numerical guards. The cost refuses an odd or non-integer power; the benchmark grid holds only the even schedules $t=2$, $t=4$, $(4,2)$ and $(4,2)$ with a design efficiency of 0.8, so no entry of the grid is refused and no benchmark row of conjugate-gradient minimization is skipped.

\textbf{Acceptance criterion and gate.} The analytic gradient agrees with finite differences, and a square lattice of flat-top spots reaches an error below one percent (Table~\ref{tab:ladder}). The gate checks the gradient for $t=2$ and $t=4$ on a shaped target and on traps by central differences in six directions, runs the lattice on a $64$-pixel SLM padded twice with 1500 iterations and two seeds, and checks that the cost decreases monotonically. The measurement of this report is: \acc{M10}. Harte et al.~\cite{harte} report 0.58 percent at 4000 iterations for the same figure.

\textbf{Why the gradient costs two transforms.} The cost of Eq.~\eqref{eq:m10} depends on the phase only through the intensity $|\tilde\psi|^2$ of the focal field, which is a linear function of $e^{i\phi}$. The derivative with respect to a pixel phase therefore needs the focal field itself, one forward transform, and the sum $X$ of Eq.~\eqref{eq:m10-x}, which is an inverse transform of the weighted residual, one adjoint transform. No derivative of the transform is evaluated numerically, and the port needs no automatic differentiation. The line search multiplies the number of cost evaluations, each a forward transform, which is why the cost of an iteration is dominated by it, and why the bake budget of a few tens of iterations leaves the method far from its limit on the trap arrays. The sequential powers for fields, the powered cost first and the basic cost after, use the second stage to remove the amplitude error that the first stage leaves while keeping the vortex-free phase that the first stage found.

\textbf{Limits.} The strong-Wolfe search is DinHot's choice. The model is the ideal device. The cost grows with the pixel count, so full-size runs are baked; and the bake budget of 30 iterations from a random start is far from the lattice test, so the trap-array rows of the bake are far from converged and are not a statement about the limit of the method. The module is \texttt{m10\_cg} on the optimization helpers of \texttt{optim/\_common}.

\subsection{Wirtinger gradient descent with Adam (M11)}\label{sec:m11}

The hologram is treated as a differentiable program: any linear forward model (the point-trap form, the Fourier form, a device) can be optimized by a first-order method if the gradient with respect to the phase is available. The solution replicated is Wirtinger holography of Chakravarthula, Peng, Kollin, Fuchs and Heide~\cite{chakra}, with the Adam optimizer of Kingma and Ba~\cite{kingma} (named in the module; the original was not re-read). It is the bridge between the iterative methods and the learned ones. The loss is minimized over the phase with a differentiable propagation $P$ and $H=c\,e^{i\phi}$; the gradient of the real loss with respect to the complex field is the Wirtinger derivative, pulled back through $P$ by the chain rule,
\begin{equation}
\frac{\partial L}{\partial\phi_j}=2\operatorname{Re}\!\big(i\,f_j\,\overline{(P^{H}G)_j}\big),\qquad G=\frac{\partial L}{\partial z^*}=L'(I)\,z,
\label{eq:m11}
\end{equation}
so that any linear forward model works. The loss is the quadratic intensity loss $C_2=\sum(I_m-T_m)^2/\sum T_m^2$ of the optimization module, with the design efficiency of conjugate-gradient minimization; the penalty term of the source has weight zero, because a phase penalty has no role for a phase-only hologram of traps. Adam updates, with bias-corrected moments,
\begin{equation}
m\leftarrow\beta_1m+(1-\beta_1)g,\quad v\leftarrow\beta_2v+(1-\beta_2)g^2,\quad x\leftarrow x-\mathrm{lr}\,\frac{\hat m}{\sqrt{\hat v}+\epsilon}.
\label{eq:m11-adam}
\end{equation}
The source uses L-BFGS for quadratic losses and Adam with a small rate for learned losses; the start is $n_{\mathrm{pix}}^2$ uniform random phases.

\textbf{Pseudocode.} Algorithm~\ref{alg:m11} lists the loop.

\begin{algorithm}[t]
\caption{Wirtinger gradient descent with Adam on the quadratic intensity loss; the history records the state before each step, so the cost may rise.}
\label{alg:m11}
\begin{algorithmic}[1]
\Require target, learning rate $\mathrm{lr}$, $\beta_1=0.9$, $\beta_2=0.999$, $\epsilon=10^{-8}$, iteration budget $K$, start $\phi^0$
\State $\phi\gets\phi^0$ (uniform random phases when no start is given); $m\gets0$; $v\gets0$
\For{$i=1,\dots,K$}
  \State forward model: $z\gets$ field of $\phi$; $L\gets C_2$; adjoint pass: $g\gets2\operatorname{Re}\big(i\,f\,\overline{P^HG}\big)$ with $G=L'(I)\,z$ (Eq.~\eqref{eq:m11}); record the state
  \State $m\gets\beta_1m+(1-\beta_1)g$; $v\gets\beta_2v+(1-\beta_2)g^2$; $\hat m\gets m/(1-\beta_1^i)$; $\hat v\gets v/(1-\beta_2^i)$
  \State $\phi\gets\phi-\mathrm{lr}\,\hat m/(\sqrt{\hat v}+\epsilon)$
\EndFor
\State \Return $\phi$ wrapped to $[0,2\pi)$
\end{algorithmic}
\end{algorithm}

\textbf{Parameters and how they were chosen.} The learning rate has the default \defv{M11}{lr}, tuned on the trap loss: the rate of the source moves a pixel by a small angle in 30 iterations, and the tuning record of the design document (tuning seeds, then separate seeds) shows the spread falling as the rate rises up to a value near the default, where a larger rate ties within one percent and a still larger one diverges; the smaller step is kept. The moment constants are the defaults of Kingma and Ba. The benchmark's tuning grid for the rate is 0.5, 1, 2 and 4, around the default of 1. The tuned row equals the default row on three of the seven geometries, and on the others the rate chosen at 30 iterations is within a factor of about 1.4; at 100 iterations, however, the tuned rows of G1, G5 and G6 are four to eight times worse than the default (a relative standard deviation of 8.5, 9.8 and 9.0 percent against 1.2, 2.2 and 1.2), so a rate tuned for a short budget can hurt a long one. A mutation run with the rate of the source fails the acceptance test, which shows that the tuned rate is what makes the row pass.

\textbf{Complexity.} One forward and one adjoint evaluation per iteration, two Fourier transforms in the Fourier form, so that the cost per iteration is the same as for the conjugate gradient without its line search; the method runs both live at small sizes and baked at full size. Adam is not a descent method, so the history records the state before each step and the cost may rise. With a device the traps must lie on bins with no depth offset and no vortex.

\textbf{Acceptance criterion and gate.} The criterion is a lower trap-intensity error than GS at equal iterations on the $10\times10$ task (Table~\ref{tab:ladder}); the gate runs three seeds of 30 iterations and requires a lower spread per seed and in median than the in-test copy of the GS rule, which was cross-checked against plain GS and gives the same numbers to four digits. Companion tests check the gradient on a set with depth offsets, a vortex and unequal weights. The measurement of this report is: \acc{M11}. Because the derivative is pulled back by the chain rule, the same code differentiates the point-trap form, the Fourier form and the device form; the only requirement on a forward model is that its adjoint is available, and the engine provides it for each. The tuning record is therefore a statement about one loss on one family of grids: a rate that is right for the $10\times10$ task is not guaranteed to be right for a crowded array or a shaped field. The method stops at a few percent in 30 iterations, since it is first order; the tuned rate is specific to this loss and these grids; and on a shaped target scored on a device the rate of the trap loss is untuned. The module is \texttt{m11\_wirtinger\_adam}.

\subsection{Closed-loop weighted GS on an emulated camera (M12)}\label{sec:m12}

Camera-in-the-loop correction measures the trap intensities and rescales the target until the traps are equal. Two published procedures are replicated, on an emulation. Nogrette et al.~\cite{nogrette} reduce a 19 percent root-mean-square trap-depth spread to 1.4 percent in 20 iterations at gain $G=0.7$, and Kim et al.~\cite{kim2019} reduce a 22 percent spread to 1.59 percent in three phase-fixed steps. The camera is an emulation of per-trap fixed gain error, additive noise, blur and roll-off, so a result shows that the loop works against those stated errors and not against an instrument; a measured hardware loop is a documented frontier of the product. One feedback iteration measures the trap intensities of the current hologram, rescales the target and runs a number of inner iterations from the current phase; the open-loop start is a number of iterations of weighted GS from the random start.

\textbf{Update rules.} With $\rho_m=(I_m/w_m)/\langle I/w\rangle$ the measured intensity relative to its target, the Nogrette variant runs plain GS inside and rescales the current target intensity,
\begin{equation}
w_m^{T}\leftarrow w_m^{T}\,\frac{1}{1-G(1-\rho_m)},
\label{eq:m12}
\end{equation}
implemented on amplitudes with a floor on $\rho$. Nogrette et al.~\cite{nogrette} compute, after each measurement, a new target intensity $I'_i=\bar I/[1-G(1-I_i/\bar I)]$ from the measured trap intensities $I_i$ alone (their Eq.~3, gain $G\simeq0.7$) and run the Gerchberg--Saxton iteration again from the previous phase. The engine applies the same factor as a multiplier of the running target intensity, Eq.~\eqref{eq:m12}; this cumulative form agrees with the source at the first correction from a uniform target, and from the second correction on it keeps the product of all earlier factors, an integral action that stops changing only at $\rho=1$, where the source recomputes the target from the latest measurement alone. Read literally, the printed rule does not close the loop: a uniform measurement returns the target to the base design, so uniformity is not its fixed point, and implemented as the variant \texttt{nogrette-printed} it ends, after 20 corrections with $G=0.7$ on seeds 1, 2 and 3, at a spread of 15.4, 16.8 and 18.2 percent, where the cumulative form ends at 0.84, 1.07 and 2.18 percent. The source reports a fall from 19 to 1.4 percent in 20 iterations, which only the cumulative reading reproduces, so DinHot reads the measured intensity of Eq.~3 as relative to the current target; the results of this method are those of the cumulative form. The Kim variant uses the square-root weights of the source,
\begin{equation}
T^{(j)}_m=\sqrt{\frac{\langle I\rangle}{I_m}}\;T^{(j-1)}_m,
\label{eq:m12-kim}
\end{equation}
at full gain, and the first inner iteration displays $\arg\mathcal F^{-1}[Te^{i\psi}]$ with the focal phase frozen at that of the current hologram and no weight correction on top: a second correction overshoots and the loop diverges, which was measured. Later inner iterations add the phase-fixed weights of phase-fixed weighted GS, reset at every step.

\textbf{Pseudocode.} Algorithm~\ref{alg:m12} lists a run; the variant that reads the update as printed, without memory, replaces $a_m\gets a_m\sqrt{\cdot}$ by $a_m\gets a^{\mathrm{base}}_m\sqrt{\cdot}$.

\begin{algorithm}[t]
\caption{Closed-loop weighted GS on the emulated camera, the two variants (the cumulative reading of the Nogrette rule and the square-root rule of Kim); the history records the spread of the measured intensities.}
\label{alg:m12}
\begin{algorithmic}[1]
\Require traps on FFT bins, camera, variant, gain $G$, inner iterations $n_{\mathrm{in}}$, start iterations $n_0$, feedback budget $K$
\State $\phi\gets$ $n_0$ iterations of weighted GS from the random superposition; $a\gets a^{\mathrm{base}}$; measure the camera readings $\mathrm{meas}$ of $\phi$
\For{$k=1,\dots,K$}
  \State $\rho_m\gets\max\{(\mathrm{meas}_m/w_m)/\langle\mathrm{meas}/w\rangle,\,10^{-3}\}$
  \If{variant is Nogrette} \State $a_m\gets a_m\sqrt{1/(1-G(1-\rho_m))}$; $n_{\mathrm{in}}$ iterations of plain GS on the target amplitudes $a$ from the current phase
  \Else \State $a_m\gets a_m/\sqrt{\rho_m}$; $g_m\gets1$; $\psi\gets$ none
    \For{$j=1,\dots,n_{\mathrm{in}}$}
      \State $V\gets$ trap fields of $\phi$; \textbf{if} $\psi$ is none \textbf{then} $\psi\gets\arg V$ \textbf{else} $b_m\gets|V_m|/a_m$; $g_m\gets g_m\langle b\rangle/b_m$ \textbf{end if}
      \State $\phi\gets\arg\mathcal F^{-1}[\sum_m g_ma_me^{i\psi_m}\delta(\mathbf u-\mathbf u_m)]$
    \EndFor
  \EndIf
  \State measure the camera readings of $\phi$; record the spread of $\mathrm{meas}/w$
\EndFor
\State \Return $\phi$
\end{algorithmic}
\end{algorithm}

\textbf{Parameters.} The variant, the gain (default \defv{M12}{gain}) and the numbers of inner and start iterations are listed in Table~\ref{tab:defaults}. The gain is that of Nogrette et al. The inner iteration count has two measured defaults: with more than one inner GS iteration the loop settles at several percent because the algorithm's own spread is reintroduced between measurements, and for the Kim variant three inner iterations reach a fraction of a percent where one does not. The default camera has a gain error of 0.19 and no roll-off, which is Nogrette's spread.

\textbf{The emulated camera.} At the bin nearest to each trap the camera returns the intensity map, optionally blurred by a Gaussian and multiplied by the pixel-aperture roll-off, times a fixed gain per trap identity $\exp(sz-s^2/2)$, with $s$ the stated gain error and $z$ the first normal of a stream seeded from the seed and the trap identity, so that the gain of a trap does not depend on the order or composition of the trap set; additive Gaussian noise is clipped at zero. Every error is a parameter with a stated meaning. The emulator was first proposed as a method and dropped from the ladder: an emulator that reproduces published feedback curves is a validation tool for the closed-loop method, not a method that answers the problem differently.

\textbf{Acceptance criterion and gate.} From the published open-loop spreads the loop reaches the published curves within a factor of 1.5 (Table~\ref{tab:ladder}). The gate runs two seeds and asserts a final measured spread within the factor and below one eighth of the open-loop spread, a plain-GS start, the Kim variant in three steps, convergence with a noisy blurred camera, and determinism. The measurement of this report is: \acc{M12}. The loop reaches the published level for the Nogrette variant and goes below it for the Kim variant, because the emulated camera has no noise. Two findings are kept as results: more than one inner GS iteration degrades the Nogrette loop, and a second weight correction in the first Kim step diverges. Only the end points of the published curves are compared, and the conventional-WGS clause of the research record was not run.

\textbf{Limits and complexity.} The method accepts traps on bins only (Fourier form, no depth offset, no vortex). One feedback iteration is the inner iterations plus one camera evaluation. The case feedback driver is a separate loop that applies the Nogrette factor to the case's weight-aware method and is the subject of Section~\ref{sec:results}. The module is \texttt{m12\_closed\_loop} on the camera of \texttt{core/camera}.
